\documentclass[11pt,a4paper]{article}
\usepackage[T1]{fontenc}
\usepackage[utf8]{inputenc}
\usepackage{lmodern}
\usepackage[margin=25mm,headheight=15pt]{geometry}
\usepackage{amsmath,amssymb,amsthm,mathtools,aliascnt}
\usepackage{microtype,booktabs,tabularx,array,longtable}
\usepackage[dvipsnames]{xcolor}
\usepackage{enumitem,fancyhdr,titlesec}
\usepackage[colorlinks=true,linkcolor=MidnightBlue,citecolor=MidnightBlue,urlcolor=MidnightBlue]{hyperref}
\usepackage[nameinlink,noabbrev]{cleveref}
\definecolor{Ink}{HTML}{21364A}
\definecolor{Accent}{HTML}{45605C}
\titleformat{\section}{\Large\bfseries\color{Ink}}{\thesection}{.65em}{}
\titleformat{\subsection}{\large\bfseries\color{Accent}}{\thesubsection}{.65em}{}
\pagestyle{fancy}\fancyhf{}
\fancyhead[L]{\small Surcomplex differential systems}
\fancyhead[R]{\small Spectra, phases, and positive metrics}
\fancyfoot[C]{\thepage}
\setlength{\parskip}{.25em}
\setlength{\parindent}{1.1em}
\setlist{itemsep=.15em,topsep=.35em,leftmargin=2em}
\setcounter{tocdepth}{2}
\emergencystretch=2em
\numberwithin{equation}{section}
\newtheorem{theorem}{Theorem}[section]
\crefname{theorem}{Theorem}{Theorems}
\crefname{section}{Section}{Sections}
\newaliascnt{proposition}{theorem}
\newtheorem{proposition}[proposition]{Proposition}
\aliascntresetthe{proposition}
\crefname{proposition}{Proposition}{Propositions}
\newaliascnt{lemma}{theorem}
\newtheorem{lemma}[lemma]{Lemma}
\aliascntresetthe{lemma}
\crefname{lemma}{Lemma}{Lemmas}
\newaliascnt{corollary}{theorem}
\newtheorem{corollary}[corollary]{Corollary}
\aliascntresetthe{corollary}
\crefname{corollary}{Corollary}{Corollarys}
\theoremstyle{definition}
\newaliascnt{definition}{theorem}
\newtheorem{definition}[definition]{Definition}
\aliascntresetthe{definition}
\crefname{definition}{Definition}{Definitions}
\newaliascnt{example}{theorem}
\newtheorem{example}[example]{Example}
\aliascntresetthe{example}
\crefname{example}{Example}{Examples}
\theoremstyle{remark}
\newaliascnt{remark}{theorem}
\newtheorem{remark}[remark]{Remark}
\aliascntresetthe{remark}
\crefname{remark}{Remark}{Remarks}
\newcommand{\No}{\mathbf{No}}
\newcommand{\K}{\mathbb K}
\newcommand{\R}{\mathbb R}
\newcommand{\C}{\mathbb C}
\newcommand{\Q}{\mathbb Q}
\newcommand{\N}{\mathbb N}
\newcommand{\Z}{\mathbb Z}
\newcommand{\T}{\mathbb T}
\newcommand{\OO}{\mathcal O}
\newcommand{\mm}{\mathfrak m}
\newcommand{\PP}{\mathcal P}
\newcommand{\II}{\mathcal I}
\newcommand{\st}{\operatorname{st}}
\newcommand{\pp}{\operatorname{pp}}
\newcommand{\supp}{\operatorname{supp}}
\newcommand{\diag}{\operatorname{diag}}
\newcommand{\tr}{\operatorname{tr}}
\newcommand{\GL}{\operatorname{GL}}
\newcommand{\SL}{\operatorname{SL}}
\newcommand{\UU}{\operatorname U}
\newcommand{\SO}{\operatorname{SO}}
\newcommand{\Herm}{\operatorname{Herm}}
\newcommand{\Sym}{\operatorname{Sym}}
\newcommand{\Hom}{\operatorname{Hom}}
\newcommand{\End}{\operatorname{End}}
\newcommand{\Specph}{\operatorname{Ph}}
\newcommand{\Span}{\operatorname{span}}
\newcommand{\Mat}{\operatorname{Mat}}
\newcommand{\Frac}{\operatorname{Frac}}
\newcommand{\dlog}{\operatorname{dlog}}
\newcommand{\abs}[1]{\left|#1\right|}
\newcommand{\norm}[1]{\left\lVert#1\right\rVert}
\newcommand{\repocommit}{aa846271b4dcae2c055b216126a87210292ec19b}
\newcolumntype{L}{>{\raggedright\arraybackslash}X}
\hypersetup{pdftitle={A Spectral--Integral Classification of Surcomplex Unitary Differential Systems},pdfauthor={Research manuscript prepared for Vladimir Reshetnikov},pdfsubject={Berarducci--Mantova derivation, spectral phase formula, invariant metrics, Ermakov--Pinney equation}}
\begin{document}
\hypersetup{pageanchor=false}
\begin{titlepage}
\centering
\vspace*{8mm}
{\large\scshape Research development for the Surreal project\par}
\vspace{12mm}
{\Huge\bfseries\color{Ink} A Spectral--Integral\par Classification of\par Surcomplex Unitary\par Differential Systems\par}
\vspace{8mm}
{\Large Purely infinite phases, non-Abelian integration,\par and positive invariant metrics\par}
\vspace{9mm}
{\large September 21, 2026\par}
\vspace{8mm}
\begin{minipage}{.95\textwidth}
\begin{abstract}
Fix the Berarducci--Mantova derivation on the surreal field, with
$\partial\omega=1$, and extend it to $\K=\No[i]$. Let
$\II=\partial\OO$ be the real coefficients with a finite primitive, and let
$\Psi(b)$ be the purely infinite part of any primitive of $b$.
We prove that, for an arbitrary finite Hermitian matrix $H$ over $\K$,
the complete differential phase spectrum of $\partial y=iHy$ is
\[
 \bigl\{\Psi(\lambda_1(H)),\ldots,\Psi(\lambda_n(H))\bigr\},
\]
where the $\lambda_j(H)$ are its ordinary algebraic eigenvalues in $\No$.
The formula needs neither commuting coefficients nor a spectral gap.
It yields an exact unitary normal form, a sharp finite-primitive
perturbation principle, and a normalized noncommutative integration theorem.

The proof combines established differential-field existence theorems with
an explicit phase classification and an order-theoretic estimate on bounded
gauges. We also classify all positive horizontal metrics of arbitrary
surcomplex systems, derive the real rotation-block form, and obtain a
positive Ermakov--Pinney existence/uniqueness dichotomy. An exact Hahn
construction for the oscillatory Airy equation illustrates the nonlinear
consequences. Further sections describe phase projectors, rational and
polynomial first integrals, and finite Picard--Vessiot tori.
\end{abstract}
\end{minipage}
\vfill
\begin{minipage}{.95\textwidth}\small
\textbf{Research status.} The proposed new contribution is the precise
spectral--integral theorem and its matrix consequences, not the construction
of the surreal derivation or the general universal exponential extension.
No claim to settle a named published open problem is made. Novelty is based
on a targeted literature search, not an exhaustive priority determination.
Proofs are supplied relative to explicitly stated published inputs; they
have not been independently refereed or proof-assistant checked.
\end{minipage}
\end{titlepage}
\hypersetup{pageanchor=true}
\pagenumbering{roman}
\tableofcontents
\clearpage\pagenumbering{arabic}

\section{The question, the contribution, and prior coverage}\label{sec:intro}

\subsection{From missing oscillations to a matrix classification}
The surcomplex field is algebraically closed, but algebraic closedness does
not make it a field containing every differential exponential. For its
intrinsic Berarducci--Mantova derivation, $\partial y=iy$ has no nonzero
solution. The obstruction is not the magnitude of the coefficient: even
$i/\omega$ is obstructed, because its imaginary part has the infinite
primitive $\log\omega$. Conversely, an infinitesimal with a finite primitive
does permit a unit-modulus solution.

The repository already develops this rank-one obstruction and several
associated oscillatory extensions. The next question is substantially
stronger: how can one recognize the missing oscillations of a \emph{coupled},
variable-coefficient matrix system? In particular, can the algebraic
spectrum of a Hermitian coefficient matrix detect them, despite the fact
that the matrix need not commute with its derivative or with any of its
primitives?

The answer proved here is yes, after passing from a real coefficient to the
purely infinite part of its primitive. This is not a claim that pointwise
matrix diagonalization commutes with differentiation. It does not. The
additional connection term is present, but its entries have finite
primitives, and that is precisely the amount of information the phase
invariant discards.

\subsection{The main statement in one display}
Write $F=\No$, $\K=F[i]$, and let $\lambda_1\le\cdots\le\lambda_n$ be the
algebraic eigenvalues of $H=H^*\in\Mat_n(\K)$. Then there is a unitary
$U\in\UU_n(\K)$ such that
\begin{equation}\label{eq:main-intro}
 U^*(iH)U-U^*\partial U
   =i\diag\bigl(\partial p_1,\ldots,\partial p_n\bigr),
 \qquad p_j=\Psi(\lambda_j).
\end{equation}
Here $\Psi(b)=\pp(B)$ for any $B\in F$ with $\partial B=b$. The multiset
$\{p_j\}$ is a complete invariant for differential unitary equivalence.
The $p_j$ are \emph{purely infinite surreal numbers}, not finite angles;
when $p_j\ne0$, the missing symbol $e^{ip_j}$ must be adjoined in a
differential extension rather than evaluated in $\K$.

Three consequences distinguish \eqref{eq:main-intro} from an asymptotic
approximation. First, perturbing every entry of $H$ by a coefficient with a
finite primitive preserves its entire differential unitary class, with no
gap assumption. Second, the dimension over $\C$ of the actual solution
space in $\K^n$ equals the number of eigenvalues whose primitives are
finite. Third, the matrix equation $\partial U=AU$ has a unique unitary
solution with $\st(U)=I$ whenever $A^*=-A$ and every entry of $A$ has a
finite primitive.

\subsection{What was inspected, and what is not claimed as new}
The repository was inspected at the tree revision
\begin{center}\small\texttt{\repocommit}.\end{center}
The documentation catalogue, the differential-equations report inventory,
and the opening portions of the spectral-theory article were the main
comparison sources~\cite{RepoMap,RepoDE,RepoSpectral}. This is a targeted
coverage audit, not a certification of every archived manuscript. The
catalogue identifies eighteen research packages and labels them as
AI-assisted drafts. The differential-equations inventory expressly leaves
general nontriangular systems and higher-rank differential modules outside
its classification. Its rank-one phase criterion, finite-primitive ideal,
and diagonal Galois tori are prior coverage, not discoveries of this paper.

The external mathematical inputs are equally important. Berarducci and
Mantova constructed the derivation~\cite{BM}. Aschenbrenner, van den Dries,
and van der Hoeven proved the elementary embedding of transseries into the
surreal differential field~\cite{ADHSur}. Their later work supplies the
linear splitting and surjectivity facts used below; it also develops
universal exponential extensions and abstract differential spectra
\cite{ADHNorm}. The group-algebra construction in \cref{sec:galois} is
therefore explicitly a concrete incarnation of existing theory.

The closest classical nonlinear antecedent is the Ermakov--Pinney
formula~\cite{Pinney}. The uniqueness conclusion below is a consequence of
the absence of oscillatory elements in the base field, rather than a new
classical integration formula. It also has an abstract antecedent in the
$\sigma$-map and Schwarz-closed-field structure discussed in
\cite[Introduction, equations (0.3)--(0.6)]{ADHNorm}.
The 2026 paper \cite{ADHSecond} develops existence and sharp uniqueness
of amplitude--phase representations over Hardy fields. That recent
antecedent is especially relevant to the second-order consequences here;
the main candidate novelty is the arbitrary-dimensional Hermitian
spectral--integral formula, not the amplitude--phase idea.

\begin{table}[ht]
\centering\small
\begin{tabularx}{\textwidth}{@{}p{.29\textwidth}L@{}}
\toprule
Component & Status in this article\\
\midrule
Surreal derivation; transfer; linear splitting and surjectivity & Published inputs, explicitly attributed.\\
Rank-one phase obstruction; universal exponential algebra & Existing repository/literature theory, reproved or specialized as needed.\\
Complete finite matrix phase reduction & An explicit development of consequences of those inputs, not a claim of a new general differential-algebra principle.\\
Spectral--integral formula for Hermitian coefficients & Principal candidate new theorem; no exact prior statement located in the targeted search.\\
Gap-free perturbation and normalized non-Abelian integration & Matrix consequences proved here; included in the candidate contribution.\\
Horizontal metric cones; real descent; Pinney dichotomy & Detailed consequences, with the classical and abstract antecedents kept visible.\\
\bottomrule
\end{tabularx}
\caption{The novelty claim is narrower than the collection of theorems proved.}
\end{table}

\subsection{Scope and reading route}
All dimensions are ordinary finite integers. All derivatives in this paper
are the fixed intrinsic derivation $\partial$, not derivatives of a
class-function of a surreal variable and not a coefficientwise derivative
with respect to an independent complex coordinate. A displayed Hahn sum
is a support-certified formal sum, not a limit in the fine order topology
of the full surreal field.

The main proof is in \cref{sec:spectralintegral}. Its ingredients are the
finite-primitive order ideal in \cref{sec:ideal}, the phase normal form in
\cref{sec:normal}, and the elementary finite-dimensional spectral estimates
in \cref{sec:hermitian}. The reader interested in the nonlinear consequence
can then proceed directly to \cref{sec:pinney,sec:airy}.

\section{The differential field and its finite-primitive ideal}\label{sec:ideal}

\subsection{Foundational inputs}
We use $F=\No$ with the Berarducci--Mantova derivation, normalized by
$\partial\omega=1$, and $\K=F[i]$. The following are background theorems,
not reconstructed here:
\begin{enumerate}
\item $F$ is real closed, $\K$ is algebraically closed, and the constant
fields of $F$ and $\K$ are $\R$ and $\C$, respectively.
\item $\partial:F\to F$ is onto, is compatible with the surreal real
exponential, and commutes with summable surreal families. Thus
$\partial\exp(f)=\exp(f)\partial f$.
\item The $H$-field positivity rule holds:
\[
 f>\R\quad\Longrightarrow\quad\partial f>0.
\]
Here $f>\R$ means that $f$ exceeds every real constant.
\item The differential field $F$ is elementarily equivalent to the
transseries field $\T$, with the natural inclusion of $\T$ into $F$
elementary in the ordered differential-field language.
\end{enumerate}
The derivation inputs are due to \cite{BM}; the last statement is
\cite[Theorem 1]{ADHSur}. The extension to $\K$ is
$\partial(a+ib)=\partial a+i\partial b$ and commutes with conjugation.

Let $\OO\subset F$ be the finite elements and $\mm\subset\OO$ the
infinitesimals. Every $f\in F$ has a unique normal-form decomposition
\begin{equation}\label{eq:decomposition}
 f=p+c+\varepsilon,
 \qquad p\in\PP,\quad c\in\R,\quad\varepsilon\in\mm.
\end{equation}
The class $\PP$ consists of zero and those normal forms supported entirely
on monomials greater than $1$. We call it the \emph{purely infinite part};
$\pp(f)=p$. Every nonzero member of $\PP$ is infinite in absolute value.
The standard part of a finite $f$ is the real coefficient $c$ in
\eqref{eq:decomposition}. Set
$\OO_{\C}=\OO+i\OO$ and $\mm_{\C}=\mm+i\mm$.

\begin{definition}
The \emph{finite-primitive ideal} is
\[
 \II:=\partial\OO=\partial\mm\subset F,
 \qquad \II_{\C}:=\II+i\II.
\]
For $b\in F$, define
\begin{equation}\label{eq:Psi}
 \Psi(b):=\pp(B)\quad\text{where }\partial B=b.
\end{equation}
\end{definition}
The equality $\partial\OO=\partial\mm$ follows by removing the standard
part. The map $\Psi$ is independent of the primitive, since two primitives
differ by a real constant. It is $\R$-linear and onto $\PP$, with kernel
$\II$. These are identities of classes and class maps, with each individual
normal-form support a set.

\subsection{The estimate that makes the matrix theorem possible}
\begin{lemma}[Convexity and bounded multiplication]\label{lem:ideal}
The class $\II$ is a convex additive subgroup of $F$, is an $\OO$-module,
and is contained in $\mm$. In particular,
\[
 a\in\II,\quad |b|\le |a|\quad\Longrightarrow\quad b\in\II,
 \qquad \OO\II\subseteq\II.
\]
\end{lemma}
\begin{proof}
Additivity is immediate. Suppose $a=\partial f$ with $f$ finite and
$|b|\le |a|$. Choose $g$ with $\partial g=b$. If $g$ were infinite, replace
$g,b$ by their negatives if necessary so that $g>\R$. Choose $h=f$ or
$h=-f$ so that $\partial h=|a|$. Then $g-h>\R$, but
\[
 \partial(g-h)=b-|a|\le0,
\]
contradicting the $H$-field rule. Therefore $g$ is finite, so $b\in\II$.
This proves the required convexity. If $c\in\OO$, choose a positive
integer $N$ with $|c|\le N$. Since $Na\in\II$ and
$|ca|\le N|a|$, convexity gives $ca\in\II$.

Finally, let $f$ be finite and let $r>0$ be real. If
$|\partial f|\ge r$, choose $s\in\{1,-1\}$ with
$s\partial f=|\partial f|$. The number $\omega-sf/r$ is positive infinite,
whereas its derivative is $1-|\partial f|/r\le0$, another contradiction.
Thus $|\partial f|<r$ for every positive real $r$.
\end{proof}

\begin{lemma}[Order of the phase quotient]\label{lem:phaseorder}
The map $\Psi:F\to\PP$ preserves order. It induces an order-preserving
$\R$-linear bijection $F/\II\to\PP$, where the quotient is the ordered
additive quotient. Also, for $p,q\in\PP$,
\[
 p<q\quad\Longleftrightarrow\quad\partial p<\partial q.
\]
\end{lemma}
\begin{proof}
If $b\ge0$ and $\partial B=b$, then $B$ cannot be negative infinite:
$-B>\R$ would imply $-\partial B>0$. Therefore $\pp(B)\ge0$, proving
order preservation. The kernel and surjectivity statements were established
above. If $p<q$, then $q-p>\R$, hence
$\partial(q-p)>0$. Injectivity of $\partial$ on $\PP$ and the total order
give the converse.
\end{proof}

\begin{remark}[Not a field quotient]
The notation $F/\II$ in \cref{lem:phaseorder} does \emph{not} denote a
quotient field or a quotient ring. The ideal property is only for the
finite-element ring $\OO$. Multiplication by an infinite element need not
preserve $\II$. The spectral theorem below takes algebraic eigenvalues
\emph{before} passing to their additive classes modulo $\II$.
\end{remark}

\begin{example}[Small is not the same as integrable]
The coefficient $\omega^{-2}$ belongs to $\II$, since it is the derivative
of $-\omega^{-1}$. The infinitesimal $\omega^{-1}$ does not: its primitives
are $\log\omega+c$. Thus $\II\subsetneq\mm$. The distinction persists at
arbitrarily slow logarithmic scales; it is not a universal single valuation
threshold on the entire surreal field.
\end{example}

\subsection{Rank-one phases, with their prior provenance}
For $a\in\K$, put
\begin{equation}\label{eq:Phi}
 \Phi(a):=\Psi(\operatorname{Im}a)\in\PP.
\end{equation}
The following rank-one statement is already present in the repository
\cite{RepoDE}. We include the proof because it is the algebraic test behind
every uniqueness argument in this paper.

\begin{proposition}[Logarithmic derivatives]\label{prop:dlog}
One has
\begin{equation}\label{eq:dlog}
 \left\{\frac{\partial z}{z}:z\in\K^\times\right\}
   =F+i\II=\ker\Phi.
\end{equation}
Consequently $\partial z=az$ has a nonzero solution in $\K$ exactly when
$\Phi(a)=0$. Every $a\in\K$ can be written
\begin{equation}\label{eq:rankone}
 a=\frac{\partial s}{s}+i\partial p,
 \qquad s\in\K^\times,\quad p=\Phi(a).
\end{equation}
\end{proposition}
\begin{proof}
For $z\ne0$, write $z=ru$ with $r=\sqrt{z\bar z}>0$ and $u\bar u=1$.
The coordinates of $u$ are finite, and its standard part is a constant
$c\in\C$ of modulus one. Hence $u/c=1+\varepsilon$ with
$\varepsilon\in\mm_{\C}$. The formal logarithm
\[
 L=\sum_{n\ge1}\frac{(-1)^{n+1}}{n}\varepsilon^n
\]
is a legitimate Hahn sum. Formal logarithm identities and $u\bar u=1$
give $\bar L=-L$, so $L=i\eta$ with $\eta\in\mm$. Strong differentiation
of this sum gives $\partial u/u=i\partial\eta$. Thus
$\partial z/z=\partial r/r+i\partial\eta\in F+i\II$.

Conversely, given $a_0+ib$ with $a_0\in F$ and $b\in\II$, choose
$A\in F$ and $\eta\in\mm$ with $\partial A=a_0$ and
$\partial\eta=b$. Then
\[
 s=\exp(A)\sum_{n\ge0}\frac{(i\eta)^n}{n!}
\]
is nonzero and has logarithmic derivative $a_0+ib$. The usual exponential
and logarithm identities here are identities of formal power series
substituted at an infinitesimal; the Neumann support principle justifies
the substitutions, and strong differentiation justifies the chain rule.
For general $a$, subtract $i\partial\Phi(a)$ and apply the established
kernel statement.
\end{proof}

A useful special case is
\begin{equation}\label{eq:phase-no-solution}
 \partial z=i\partial(p-q)z,
 \quad p,q\in\PP,
 \qquad
 z\ne0\Longrightarrow p=q\text{ and }z\in\C^\times.
\end{equation}
This does not identify $p$ with a finite argument of a unit complex number.
It identifies the \emph{obstructed accumulated phase} of a differential
equation.

\section{The linear existence input and its precise transfer}\label{sec:existence}

\subsection{Two properties, not a homogeneous-solution theorem}
We require the following established properties of $\T[i]$:
\begin{equation}\label{eq:inputs}
 \begin{split}
 &\text{every positive-order monic scalar differential operator splits into
 first-order factors;}\\
 &\text{every nonzero scalar linear differential operator is onto.}
 \end{split}
\end{equation}
They are recorded explicitly in \cite[Lemma 2.4.19]{ADHNorm}; the
surjectivity needed for our diagonal-reduction proof is only the order-one
case. In that source, ``linearly closed'' means the first property, not the
assertion that every homogeneous equation has a full set of solutions in
the base field. Confusing these meanings would make $\partial y=iy$
appear solvable and would destroy the argument.

\begin{proposition}[Transfer to surcomplex coefficients]\label{prop:transfer}
Both properties in \eqref{eq:inputs} hold over $\K=\No[i]$. In particular,
for every $a,f\in\K$ there is $u\in\K$ with
\[
 \partial u-au=f.
\]
\end{proposition}
\begin{proof}
For each fixed order $n$, the assertion that every monic operator of order
$n$ factors into $n$ first-order operators is a first-order sentence in the
differential-field language. Expanding a product uses the finite rule
$\partial a=a\partial+\partial(a)$, so its coefficient identities are
polynomial identities in finitely many coefficients and their derivatives.
Likewise, order-$n$ surjectivity is a first-order sentence, with a nonzero
leading coefficient as a hypothesis.

Interpret a complex coefficient as an ordered pair of real coefficients;
complex multiplication, addition, and differentiation are definable in this
interpretation. The corresponding sentences about $\T[i]$ are therefore
sentences about the real differential field $\T$. Elementary equivalence
with $\No$ transfers these \emph{universally quantified} sentences, giving
the claims for arbitrary surreal coefficients, not merely coefficients in
the embedded copy of $\T$.
\end{proof}

No effective factorization algorithm for arbitrary surreal input follows
from this transfer. It supplies finite witnesses for each finite system.
This logical distinction is important: later exact formulas reduce the
phase problem to algebraic eigenvalues and scalar primitives, but neither
operation is asserted to be computable on unrestricted normal forms.

\subsection{A cyclic-vector proof for the finite algebra used below}
We give the needed cyclic-vector argument to make the passage from scalar
operators to matrix systems explicit. A differential module is a finite
free $\K$-space $M$ with an additive operator $\nabla$ satisfying
$\nabla(fv)=\partial f\,v+f\nabla v$.

\begin{lemma}[Cyclic vector]\label{lem:cyclic}
Every differential module of dimension $n\ge1$ over $\K$ has a vector $v$
for which $v,\nabla v,\ldots,\nabla^{n-1}v$ is a basis.
\end{lemma}
\begin{proof}
Fix a basis $e_0,\ldots,e_{n-1}$. Extend $\nabla$ to $M[T]$ by giving $T$
derivative $1$. Recursively choose $v_0,\ldots,v_{n-1}$ so that
\[
 v_0=e_0,\qquad
 v_k=e_k-\sum_{j=0}^{k-1}\binom{k}{j}\nabla^{k-j}v_j.
\]
For $v(T)=\sum_{j=0}^{n-1}T^jv_j/j!$, the Leibniz rule gives
$\nabla^k v(T)|_{T=0}=e_k$ for $0\le k<n$. The determinant of the $n$
columns $v(T),\nabla v(T),\ldots,\nabla^{n-1}v(T)$ is consequently a
nonzero polynomial in $T$, with constant term $1$.

Substitution $T=\omega-c$, for $c\in\R$, respects differentiation. A
nonzero polynomial has only finitely many roots in a field, whereas the
values $\omega-c$ are distinct. Choose a real $c$ outside that finite
exceptional set. The specialized vector is cyclic.
\end{proof}

\begin{lemma}[A rank-one filtration]\label{lem:flag}
Every finite differential module over $\K$ has a filtration by differential
submodules whose successive quotients have dimension one.
\end{lemma}
\begin{proof}
For a cyclic vector, the module is
$\K[\partial]/\K[\partial]L$, where $L$ is its monic minimal annihilator
of order $n$. By \cref{prop:transfer}, write
$L=(\partial-b_n)R$, where $R$ is a product of $n-1$ first-order factors.
The class of $R$ is nonzero, spans a one-dimensional differential submodule,
and satisfies $\partial R=b_nR$ modulo $L$. The quotient is the module
with annihilator $R$. Induction gives the filtration. The dimension claims
follow directly from division by a monic operator: the remainder has order
strictly below that operator's order.
\end{proof}

\section{Complete phase reduction of arbitrary finite systems}\label{sec:normal}

\subsection{Gauge transformations and the normal form}
A system has the form $\partial y=Ay$, with $A\in\Mat_n(\K)$. Under
$y=Pz$, where $P\in\GL_n(\K)$, its coefficient matrix becomes
\begin{equation}\label{eq:gauge}
 A^P=P^{-1}AP-P^{-1}\partial P.
\end{equation}
The derivative term is essential. Similarity of matrices alone is not
differential equivalence.

\begin{theorem}[Complete phase normal form]\label{thm:normal}
For every $A\in\Mat_n(\K)$, there exist $P\in\GL_n(\K)$ and
$p_1,\ldots,p_n\in\PP$ such that
\begin{equation}\label{eq:normal}
 A^P=D_p:=i\diag(\partial p_1,\ldots,\partial p_n).
\end{equation}
The multiset $\Specph(A)=\{p_1,\ldots,p_n\}$ is independent of $P$ and is
a complete invariant of differential equivalence.
\end{theorem}
\begin{proof}
Apply \cref{lem:flag} to the connection $\nabla=\partial-A$ on $\K^n$.
An adapted basis makes the system upper triangular. We now split its
extensions explicitly. Suppose induction has diagonalized the first
$n-1$ coordinates, leaving
\[
 A=\begin{pmatrix}D&b\\0&a\end{pmatrix},\qquad
 D=\diag(d_1,\ldots,d_{n-1}).
\]
For $S=\left(\begin{smallmatrix}I&u\\0&1\end{smallmatrix}\right)$,
the off-diagonal column of $A^S$ is
\[
 b+Du-ua-\partial u.
\]
Each equation $\partial u_j-(d_j-a)u_j=b_j$ has a solution by
\cref{prop:transfer}. This kills the column and proves diagonalizability
by induction. Repeated diagonal entries cause no obstruction: in that case
the relevant equation is ordinary intrinsic integration.

For each resulting scalar coefficient $a_j$, use \eqref{eq:rankone} to
write $a_j=\partial s_j/s_j+i\partial p_j$. The diagonal gauge with entries
$s_j$ gives \eqref{eq:normal}.

For uniqueness, suppose two normal forms $D_p,D_q$ are related by an
invertible matrix $T$. Then
\begin{equation}\label{eq:intertwiner}
 \partial T=D_pT-TD_q,
 \qquad
 \partial T_{jk}=i\partial(p_j-q_k)T_{jk}.
\end{equation}
By \eqref{eq:phase-no-solution}, an entry can be nonzero only when
$p_j=q_k$, and then it is constant. Invertibility forces equality of each
phase multiplicity. Conversely, equal phase multisets give the same normal
form after permuting coordinates, proving completeness.
\end{proof}

\subsection{Solution spaces, projectors, and resonances}
For $p\in\PP$, let $m_p(A)$ be its multiplicity, with omitted phases having
multiplicity zero.

\begin{corollary}[All inhomogeneous systems are solvable]\label{cor:inhomogeneous}
For every $A\in\Mat_n(\K)$ and $f\in\K^n$, the equation
$\partial y-Ay=f$ has a solution in $\K^n$. Its solution set is an affine
space over $\C$ of dimension $m_0(A)$. In particular,
\[
 \dim_{\C}\ker_{\K^n}(\partial-A)=m_0(A).
\]
\end{corollary}
\begin{proof}
Use the gauge in \cref{thm:normal} and solve each inhomogeneous scalar
equation by \cref{prop:transfer}. In the homogeneous equation a nonzero
coordinate is possible exactly for phase zero, where it is a complex
constant. The difference of two particular solutions is homogeneous.
\end{proof}

This strengthens an important limitation recorded in the repository:
outside the rank-one admissible locus, there was only a uniqueness
statement for general forcing. Here existence for \emph{every} forcing
follows from the separate published linear-surjectivity input. It does not
follow merely from the rank-one homogeneous criterion.

\begin{proposition}[Canonical phase subspaces and endomorphisms]\label{prop:projectors}
Let $E_p$ be the coordinate projector onto the phase-$p$ block of a normal
form. Then
\[
 \Pi_p:=PE_pP^{-1}
\]
is independent of the choice of normalizing $P$. The projectors satisfy
\[
 \Pi_p\Pi_q=\delta_{pq}\Pi_p,\qquad
 \sum_p\Pi_p=I,\qquad
 \partial\Pi_p=A\Pi_p-\Pi_pA.
\]
The algebra of horizontal endomorphisms of the system is isomorphic to
$\prod_p\Mat_{m_p(A)}(\C)$.
\end{proposition}
\begin{proof}
Every change between normalizing gauges is constant and block diagonal by
\eqref{eq:intertwiner}; it therefore commutes with each $E_p$. The displayed
identities follow by multiplication and differentiation. The same entrywise
calculation identifies the entire horizontal endomorphism algebra.
\end{proof}

More generally,
\[
 \dim_{\C}\Hom_{\partial}(A,B)=\sum_p m_p(A)m_p(B).
\]
There are no nonsplit extensions between finite differential modules over
$\K$, as the explicit off-diagonal elimination proves. Direct sums unite
phase multisets; tensor products add phases pairwise; duals negate phases.
Thus the finite-module theory is equivalent, in this concrete sense, to
finite-support $\PP$-graded finite-dimensional complex linear algebra.
The classification is not the assertion that every grade has a horizontal
vector in the original field.

\begin{corollary}[The determinant sees only the phase sum]\label{cor:trace}
For every $A$,
\[
 \Phi(\tr A)=\sum_p m_p(A)p.
\]
\end{corollary}
\begin{proof}
Take traces in $A=P'P^{-1}+PD_pP^{-1}$ and use the finite Jacobi identity
$\tr(P'P^{-1})=(\det P)'/\det P$, whose phase is zero.
\end{proof}
The trace cannot recover the individual phases: $\{p,-p\}$ and $\{0,0\}$
have the same sum. The Hermitian theorem below recovers the full multiset.

\section{Horizontal metrics and exact unitary reduction}\label{sec:metrics}

\subsection{Every system has a positive invariant metric}
A Hermitian matrix $G\in\Mat_n(\K)$ is positive definite when
$v^*Gv>0$ for every nonzero $v\in\K^n$. It is \emph{horizontal} for the
system $A$ when
\begin{equation}\label{eq:metric}
 \partial G+A^*G+GA=0.
\end{equation}
Formally, \eqref{eq:metric} says that the quadratic expression $y^*Gy$ is
constant along a solution, in any differential extension carrying the
compatible conjugation. It does not assert that the system already has
nonzero number-valued solutions in $\K^n$.

\begin{theorem}[Complete metric cone]\label{thm:metric}
Let $P$ normalize $A$ as in \cref{thm:normal}. Every horizontal Hermitian
matrix is uniquely of the form
\begin{equation}\label{eq:metriccone}
 G=P^{-*}\left(\bigoplus_p C_p\right)P^{-1},
 \qquad C_p\in\Herm_{m_p(A)}(\C).
\end{equation}
It is positive definite exactly when every $C_p$ is positive definite.
In particular, every finite system has a positive horizontal metric, and
the real dimension of its space of horizontal Hermitian forms is
$\sum_p m_p(A)^2$.
\end{theorem}
\begin{proof}
Put $B=P^*GP$. Since $P'=AP-PD_p$, differentiating gives
\[
 B'=-D_p^*B-BD_p=D_pB-BD_p.
\]
Thus $B_{jk}'=i\partial(p_j-p_k)B_{jk}$. Equation
\eqref{eq:phase-no-solution} makes every off-phase entry zero and every
within-phase entry constant. Hermitian symmetry gives exactly the matrices
$C_p$ in \eqref{eq:metriccone}. Congruence by an invertible matrix preserves
positive definiteness. For a constant Hermitian matrix, positivity over
$\K$ is equivalent to ordinary complex positive definiteness: diagonalize
it by an ordinary constant unitary matrix and inspect its real eigenvalues.
Taking every $C_p=I$ proves existence. Finally,
$\dim_{\R}\Herm_m(\C)=m^2$.
\end{proof}

\begin{example}[A positive metric is not an analytic stability claim]
For
\[
 A=\begin{pmatrix}1&1\\0&1\end{pmatrix},\qquad
 P=e^{\omega}\begin{pmatrix}1&\omega\\0&1\end{pmatrix},
\]
one has $P'=AP$, so both phases are zero. The metric
\[
 G=e^{-2\omega}\begin{pmatrix}1&-\omega\\-\omega&\omega^2+1\end{pmatrix}
\]
satisfies \eqref{eq:metric} and has positive determinant $e^{-4\omega}$.
The ordinary differential equation with the same printed coefficients has
growing solutions. There is no contradiction: the invariant metric itself
has a strongly varying scale. Our theorem has no uniform-equivalence-to-
the-Euclidean-norm hypothesis and makes no Lyapunov-stability assertion.
\end{example}

\subsection{When the original metric is already the identity}
\begin{theorem}[Unitary phase normal form]\label{thm:unitary}
If $A^*=-A$, its phase-normalizing matrix can be chosen unitary. For a fixed
ordering of its phases, any two such unitary normalizing matrices differ
by a constant block unitary matrix, with one block for each repeated phase.
The canonical phase projectors in \cref{prop:projectors} are then
Hermitian orthogonal projectors.
\end{theorem}
\begin{proof}
The identity matrix satisfies \eqref{eq:metric}. For any normalizing $P$,
\cref{thm:metric} therefore says that $P^*P=C$ is a constant positive
definite block Hermitian matrix. Let $R=C^{-1/2}$, computed in the ordinary
finite-dimensional complex block matrices. Then $R'=0$ and $R$ commutes
with $D_p$, so $U=PR$ is still a normalizing gauge and $U^*U=I$.
The uniqueness statement follows from \eqref{eq:intertwiner}; unitarity
makes its constant diagonal blocks unitary. Finally $UE_pU^*$ is a
Hermitian orthogonal projector.
\end{proof}

The positive square root used in this proof is of a \emph{constant complex
matrix}. There is no appeal to analytic functional calculus on a
class-sized surreal Hilbert space. More importantly, the theorem concerns
differential diagonalization: the term $U^*U'$ remains in the formula.

\section{Finite Hermitian algebra and the bounded-gauge estimate}\label{sec:hermitian}

\subsection{Spectral algebra without compactness}
We record the exact spectral estimate needed for the main theorem. It
holds over any real closed field and its complexification; thus it does
not rely on the specific surreal derivation. Related finite spectral
material is already developed in the repository~\cite{RepoSpectral}.

\begin{lemma}[Finite spectral theorem]\label{lem:spectral}
Every Hermitian matrix over $\K$ has a unitary eigenbasis and real
 eigenvalues. If $\lambda_1(H)\le\cdots\le\lambda_n(H)$, then
\begin{equation}\label{eq:minmax}
 \lambda_k(H)=\min_{\dim S=k}\ \max_{\substack{v\in S\\v^*v=1}}v^*Hv.
\end{equation}
Both extremal values in this formula are attained.
\end{lemma}
\begin{proof}
Algebraic closedness gives an eigenvector $v\ne0$ and an eigenvalue
$\lambda$. Since $v^*v>0$ and $v^*Hv$ is real, the identity
$v^*Hv=\lambda v^*v$ implies $\lambda\in F$. Normalize $v$ by the positive
square root of $v^*v$. Its orthogonal complement is invariant under $H$;
finite induction gives a unitary eigenbasis.

For any fixed $S$, the compression of $H$ to an orthonormal basis of $S$
is again a finite Hermitian matrix, so its largest eigenvalue attains the
inner maximum. Choosing $S$ spanned by the first $k$ ordered eigenvectors
gives the value $\lambda_k(H)$. Every $k$-dimensional $S$ meets the span
of the last $n-k+1$ eigenvectors nontrivially, and a unit vector in that
intersection has Rayleigh quotient at least $\lambda_k(H)$. This proves
\eqref{eq:minmax}, including attainment, by finite algebra alone.
\end{proof}

\begin{lemma}[An entrywise Weyl bound]\label{lem:weyl}
For Hermitian $H,E\in\Mat_n(\K)$, put
$m=\max_{j,k}|E_{jk}|$. Then
\begin{equation}\label{eq:weyl}
 |\lambda_j(H+E)-\lambda_j(H)|\le nm
 \qquad(1\le j\le n).
\end{equation}
If every entry of $E$ belongs to $\II_{\C}$, then every eigenvalue
difference in \eqref{eq:weyl} belongs to $\II$.
\end{lemma}
\begin{proof}
For $v^*v=1$, the triangle and Cauchy--Schwarz inequalities over the real
closed field give
\[
 |v^*Ev|\le m\left(\sum_j|v_j|\right)^2\le nm.
\]
Use the span of the first $j$ eigenvectors of $H$ in \eqref{eq:minmax} to
obtain the upper bound; interchange $H$ and $H+E$ for the lower bound.
If $E_{jk}=a+ib$ with $a,b\in\II$, then
$|E_{jk}|\le|a|+|b|\in\II$, so convexity gives $|E_{jk}|\in\II$.
The finite maximum and its multiple $nm$ belong to $\II$, and another
application of convexity finishes the proof.
\end{proof}

\subsection{Bounded gauges have integrable connection terms}
\begin{lemma}[Bounded-gauge estimate]\label{lem:boundedgauge}
If $U\in\UU_n(\K)$, every entry of $U$ and $U^*$ is finite, every entry of
$U'$ belongs to $\II_{\C}$, and every entry of $U^*U'$ belongs to
$\II_{\C}$. Moreover $U^*U'$ is skew-Hermitian.
\end{lemma}
\begin{proof}
A unitary column has sum of squared moduli $1$, so every entry has modulus
at most $1$. Differentiating a finite real or imaginary coordinate gives an
element of $\II$. Since $\II$ is an $\OO$-module, finite sums of products
of entries of $U^*$ with entries of $U'$ stay in $\II_{\C}$.
Differentiation of $U^*U=I$ gives $(U')^*U+U^*U'=0$.
\end{proof}

It is the conclusion $U^*U'\in\Mat_n(\II_{\C})$, not just its
infinitesimal size, that matters. An arbitrary infinitesimal perturbation
may accumulate an infinite phase. The bounded-primitive conclusion is what
allows eigenvalue errors to disappear under $\Psi$.

\section{The spectral--integral theorem}\label{sec:spectralintegral}

\subsection{Complete phases from algebraic eigenvalues}
\begin{theorem}[Spectral--integral classification]\label{thm:main}
Let $H=H^*\in\Mat_n(\K)$, and let its algebraic eigenvalues be ordered
$\lambda_1\le\cdots\le\lambda_n$. Define
\[
 p_j:=\Psi(\lambda_j)\in\PP.
\]
Then $p_1\le\cdots\le p_n$ and
\begin{equation}\label{eq:main}
 \Specph(iH)=\{p_1,\ldots,p_n\}.
\end{equation}
There exists $U\in\UU_n(\K)$ with
\begin{equation}\label{eq:main-gauge}
 U^*(iH)U-U^*U'=i\diag(p_1',\ldots,p_n').
\end{equation}
Two Hermitian coefficients $H_1,H_2$ give differentially unitarily
equivalent systems $y'=iH_1y$ and $y'=iH_2y$ if and only if their ordered
algebraic eigenvalues have identical images under $\Psi$, counted with
multiplicity.
\end{theorem}
\begin{proof}
By \cref{thm:unitary}, choose a unitary normalizing gauge with phases
$q_1\le\cdots\le q_n$. Rearranging its defining identity gives
\begin{equation}\label{eq:spectral-proof}
 U^*HU=\diag(q_1',\ldots,q_n')-iU^*U'.
\end{equation}
By \cref{lem:boundedgauge}, $E=-iU^*U'$ is Hermitian and has all entries in
$\II_{\C}$. By \cref{lem:phaseorder}, the derivatives $q_j'$ are in the
same order as the phases $q_j$. The entrywise Weyl bound therefore gives
\[
 \lambda_j(H)-q_j'\in\II.
\]
Apply $\Psi$ and use $\Psi(q_j')=q_j$ to obtain
$\Psi(\lambda_j(H))=q_j$. This proves \eqref{eq:main} and
\eqref{eq:main-gauge}. The order of the $p_j$ follows from order
preservation of $\Psi$. Completeness follows from \cref{thm:normal}; when
both systems are skew-Hermitian their normalizing gauges are unitary, so
the resulting equivalence is unitary as well.
\end{proof}

The formula is exact for arbitrary surreal scales and arbitrary ordinary
finite multiplicities. There is no requirement that the $\lambda_j$ be
distinct, that the nonzero gaps dominate a derivative, or that $H$ commute
with anything. At the same time, the theorem does \emph{not} remove the
connection term in \eqref{eq:spectral-proof}. That term is merely invisible
in the ordered additive quotient by $\II$.

\subsection{Gap-free perturbations and actual number-valued solutions}
\begin{corollary}[Finite-primitive perturbation invariance]\label{cor:perturbation}
If $H_1,H_2$ are Hermitian and every entry of $H_1-H_2$ lies in
$\II_{\C}$, then $y'=iH_1y$ and $y'=iH_2y$ are differentially unitarily
equivalent. This conclusion has no spectral-gap hypothesis.
\end{corollary}
\begin{proof}
By \cref{lem:weyl}, corresponding ordered eigenvalues differ by elements
of $\II$. Their $\Psi$-images are equal, so apply \cref{thm:main}.
\end{proof}

\begin{corollary}[Exact kernel count]\label{cor:kernel}
For Hermitian $H$,
\begin{equation}\label{eq:kernel-count}
 \dim_{\C}\{y\in\K^n:y'=iHy\}
   =\#\{j:\lambda_j(H)\in\II\}.
\end{equation}
The count includes algebraic multiplicity. The inhomogeneous equation
$y'-iHy=f$ is solvable for every $f\in\K^n$, and is uniquely solvable
exactly when none of the eigenvalues belongs to $\II$.
\end{corollary}
\begin{proof}
Combine \cref{thm:main,cor:inhomogeneous} with $\ker\Psi=\II$.
\end{proof}

For instance, if $H=\diag(\omega^{-2},\omega^{-1},1)$, the kernel has
complex dimension one. Replacing $H$ by any Hermitian matrix differing
from it entrywise by $\II_{\C}$ preserves that answer, even though such a
perturbation can be much larger than a gap at a more remote infinitesimal
scale.

\begin{proposition}[Sharpness of the perturbation class]\label{prop:sharp}
For scalar real coefficients $h,k$, the equations $y'=ihy$ and $y'=iky$
are differentially unitarily equivalent exactly when $h-k\in\II$.
Consequently no strictly larger additive class of scalar perturbations can
preserve the phase of every Hermitian system.
\end{proposition}
\begin{proof}
Scalar gauge equivalence is $i(h-k)=u'/u$ with $|u|=1$.
\Cref{prop:dlog} proves necessity and constructs a unit solution for
sufficiency. If $b\notin\II$, the one-dimensional systems with
coefficients $0$ and $b$ have different phases.
\end{proof}
In particular, ``infinitesimal perturbation'' cannot replace
``finite-primitive perturbation'': the scalar change $0\mapsto\omega^{-1}$
already changes the phase from $0$ to $\log\omega$.

\subsection{Why the Hermitian hypothesis is essential}
\begin{example}[A nonnormal counterexample]\label{ex:nonnormal}
Let $x=\omega$ and set
\[
 P=\begin{pmatrix}1-x^2&x\\-x&1\end{pmatrix},\qquad
 A=\begin{pmatrix}-x&1+x^2\\-1&x\end{pmatrix}.
\]
Then $\det P=1$ and $P'=AP$. Hence $A$ has phase multiset $\{0,0\}$ and a
fundamental matrix in the original real field. But
\[
 \det(zI-A)=z^2+1,
\]
so its instantaneous algebraic eigenvalues are $i$ and $-i$.
Taking primitives of their imaginary parts would incorrectly predict
$\{\omega,-\omega\}$. The similarity gauges needed for a general
nonnormal matrix need not be bounded. Their derivative correction can
therefore carry a nonzero phase. This example also rules out extending the
main formula to arbitrary matrices by simply extracting the imaginary
parts of their eigenvalues.
\end{example}

\subsection{A genuinely coupled unitary example}\label{subsec:unitaryexample}
Let $x=\omega$ and
\[
 a=\frac{1-x^2}{1+x^2},\qquad b=\frac{2x}{1+x^2},\qquad
 U=\begin{pmatrix}a&-b\\b&a\end{pmatrix},\qquad
 D=i\diag(x,x^{-1}).
\]
Define $A=U'U^{\mathsf T}+UDU^{\mathsf T}$. Then $A^*=-A$ and its exact
phase normal form is $D$, with
\[
 \Specph(A)=\left\{\frac{\omega^2}{2},\log\omega\right\}.
\]
For clarity, its entries are
\[
 A=\begin{pmatrix}
 i(a^2x+b^2/x)&-\dfrac{2}{1+x^2}+iab(x-x^{-1})\\[4pt]
 \dfrac{2}{1+x^2}+iab(x-x^{-1})&i(b^2x+a^2/x)
 \end{pmatrix}.
\]
It is neither diagonal nor a constant-matrix system. If $H=-iA$, then
\[
 \tr H=x+x^{-1},\qquad
 \det H=1-\frac{4}{(1+x^2)^2}.
\]
Thus the instantaneous eigenvalues of $H$ are \emph{not} exactly $x$ and
$x^{-1}$. Their discrepancies have finite primitives, however, so the main
theorem recovers the same two phases. Both are nonzero, and every forced
system $y'-Ay=f$ has a unique surcomplex solution. The phases are also
$\Q$-linearly independent, a fact used by the Galois interpretation below.

\section{Normalized non-Abelian integration}\label{sec:nonabelian}

\subsection{A matrix criterion with no ordering prescription}
Let
\[
 \mathfrak u_n(\II_{\C})
   =\{A\in\Mat_n(\II_{\C}):A^*=-A\},
\]
and let $\UU_n(1+\mm_{\C})$ denote the unitary matrices $U$ with
$U-I\in\Mat_n(\mm_{\C})$. This notation describes a congruence subgroup,
not an entrywise restriction that off-diagonal entries be close to $1$.

\begin{theorem}[Normalized compact integration]\label{thm:nonabelian}
The logarithmic derivative map
\begin{equation}\label{eq:nonabelian}
 \UU_n(1+\mm_{\C})\longrightarrow\mathfrak u_n(\II_{\C}),
 \qquad U\longmapsto U'U^*
\end{equation}
is a bijection of classes. Equivalently, for a skew-Hermitian $A$, the
following are equivalent:
\begin{enumerate}
\item every entry of $A$ has finite real and imaginary primitives;
\item $y'=Ay$ has a fundamental matrix in $\K$;
\item $U'=AU$ has a unitary solution with $\st(U)=I$.
\end{enumerate}
The solution in (3) is unique.
\end{theorem}
\begin{proof}
The bounded-gauge estimate shows that \eqref{eq:nonabelian} is well-defined
and that (3) implies (1). Suppose (1). For $H=-iA$, every entry belongs to
$\II_{\C}$, so \cref{lem:weyl}, applied against the zero matrix, places
every eigenvalue in $\II$. The main theorem gives a unitary $V$ with
$A^V=0$, hence $V'=AV$. Taking standard parts of $V^*V=I$ shows that
$C=\st(V)$ is an ordinary unitary matrix. The matrix $U=VC^{-1}$ has the
required derivative and standard part.

If $U_1,U_2$ are two normalized solutions, then
$(U_1^{-1}U_2)'=0$. This ratio is a constant complex matrix and has standard
part $I$, so it equals $I$. This proves uniqueness and injectivity of
\eqref{eq:nonabelian}. Clearly (3) implies (2). Conversely, a fundamental
matrix makes all phases zero by \cref{cor:inhomogeneous}; for a
skew-Hermitian system, \cref{thm:unitary} then supplies a unitary
fundamental matrix, proving (2) implies (3).
\end{proof}

This is a noncommutative existence theorem, not a formula
$U=\exp(\int A)$. Its normalization is residue data, $\st(U)=I$, not an
initial value at a real or surreal time. There is no chosen contour and
no claimed topological convergence of Picard iteration.

For the unique solution $U_A$, the ordinary product rule yields
\begin{equation}\label{eq:cocycle}
 (U_AU_B)'(U_AU_B)^*
    =A+U_ABU_A^*.
\end{equation}
The right side again lies in $\mathfrak u_n(\II_{\C})$. Thus the bijection
transports the group law of the near-identity unitary group to
$A\star B=A+U_ABU_A^*$. This is not ordinary addition unless the relevant
matrices commute.

\subsection{An exact noncommuting Hahn example}
Let $t=\omega^{-1}$, so $\partial t=-t^2$, and take the constant
skew-Hermitian matrices
\[
 X=\begin{pmatrix}0&i\\i&0\end{pmatrix},\qquad
 Y=\begin{pmatrix}i&0\\0&-i\end{pmatrix}.
\]
For $A=-t^2X-2t^3Y$, the normalized solution can be written as a Hahn series
$U=\sum_{n\ge0}U_nt^n$. The exact coefficient recurrence is
\begin{equation}\label{eq:noncommrec}
 U_0=I,\qquad
 (n+1)U_{n+1}=XU_n+2YU_{n-1},\qquad U_{-1}=0.
\end{equation}
Its first terms are
\begin{align*}
 U={}&I+tX+t^2\left(\frac{X^2}{2}+Y\right)\\
 &+t^3\left(\frac{X^3}{6}+\frac{XY}{3}+\frac{2YX}{3}\right)+\cdots.
\end{align*}
This differs from $\exp(tX+t^2Y)$ at order $t^3$ by
$(YX-XY)/6\ne0$. The nonnegative integer support makes the series strongly
summable, so the recurrence proves the equation exactly. Since
$(U^*U)'=0$ and its standard part is $I$, it is unitary and is the unique
solution supplied by \cref{thm:nonabelian}. The example provides an explicit
computation in a restricted Hahn workspace; the general theorem does not
assume that such a coefficient recurrence is effectively available.

\section{Real systems, rotation blocks, and metric signatures}\label{sec:real}

\subsection{Descent from conjugate phases}
Set
\[
 J=\begin{pmatrix}0&1\\-1&0\end{pmatrix}.
\]
For a real system, complex conjugation sends phase $p$ to phase $-p$.
The two phase multiplicities are therefore equal.

\begin{theorem}[Real rotation-block classification]\label{thm:real}
For every $A\in\Mat_n(F)$, there is $P\in\GL_n(F)$ for which
\begin{equation}\label{eq:realnormal}
 A^P=0_{m_0}\ \oplus\!\bigoplus_{p>0}
          \bigl(p'J\bigr)^{\oplus m_p},
 \qquad n=m_0+2\sum_{p>0}m_p.
\end{equation}
The positive phases $p\in\PP$ and multiplicities are unique. Every finite
real differential module is a direct sum of trivial one-dimensional
modules and nontrivial two-dimensional rotation modules. If $A$ is
skew-symmetric, $P$ can be chosen orthogonal.
\end{theorem}
\begin{proof}
The uniqueness of complex phase subspaces makes conjugation interchange
the $p$ and $-p$ subspaces. On the zero-phase subspace, its horizontal
complex vector space is stable under conjugation. Any vector is the sum of
a real vector and $i$ times a real vector; a real basis of the fixed part
is therefore a complex basis after complexification.

For $p>0$, choose normalizing columns $v_1,\ldots,v_{m_p}$ for phase $p$;
their conjugates give a basis for phase $-p$. Write $v_j=a_j+ib_j$. Since
$v_j'=Av_j-ip'v_j$,
\[
 a_j'=Aa_j+p'b_j,\qquad b_j'=Ab_j-p'a_j.
\]
The real columns $a_j,b_j$ thus give the block $p'J$ in
\eqref{eq:realnormal}. They are independent because the change between
$(v_j,\bar v_j)$ and $(a_j,b_j)$ is invertible over $\C$.
Uniqueness follows from the complex theorem. The only invariant real
subspaces in one nonzero phase pair are $0$ and the entire pair, because
a real subspace complexifies to a conjugation-stable sum of its two
one-dimensional distinct phase lines.

For skew-symmetric $A$, the identity metric is horizontal. In the real
normal form, $P^{\mathsf T}P$ is a positive constant block matrix commuting
with the rotation blocks, as is also shown explicitly below. Its positive
constant square root commutes with those blocks. Multiplication by its
inverse square root makes $P$ orthogonal without changing the normal form.
\end{proof}

\subsection{The complete real cone}
\begin{theorem}[Real metric cone and parity]\label{thm:realmetrics}
Every real system has a positive symmetric horizontal metric. For the
multiplicities in \eqref{eq:realnormal}, its cone of such metrics is
isomorphic to
\begin{equation}\label{eq:realmetriccone}
 \Sym_{m_0}^{+}(\R)\ \times\!\prod_{p>0}\Herm_{m_p}^{+}(\C).
\end{equation}
Here and below a product over an empty set is a point. The real vector
space of all horizontal symmetric forms has dimension
\[
 \frac{m_0(m_0+1)}2+\sum_{p>0}m_p^2.
\]
On each nonzero real phase block, the positive and negative indices of any
nondegenerate horizontal symmetric form are both even.
\end{theorem}
\begin{proof}
Pull a real horizontal form back by the real gauge $P$ in
\cref{thm:real}. Complex phase separation eliminates all terms between
different positive phase pairs or between such a pair and phase zero.
The zero block is an arbitrary constant real symmetric matrix.

Group a repeated positive phase as $p'J_m$, where
$J_m=\left(\begin{smallmatrix}0&I_m\\-I_m&0\end{smallmatrix}\right)$.
The remaining real symmetric horizontal matrices are constant matrices
commuting with $J_m$, hence exactly
\[
 C=\begin{pmatrix}S&T\\-T&S\end{pmatrix},
 \qquad S^{\mathsf T}=S,\quad T^{\mathsf T}=-T.
\]
This is the realification of the Hermitian matrix $S-iT$. A constant
complex unitary diagonalization of $S-iT$ becomes a real orthogonal
diagonalization of $C$, with every eigenvalue repeated twice. This proves
positivity, the cone formula, and the parity statement. Congruence by
$P^{-1}$ gives all metrics in the original coordinates.
\end{proof}

The products in \eqref{eq:realmetriccone} refer to finite collections of
constant-matrix cones; there is no integration over phases. In particular,
a real irreducible two-dimensional system has a one-dimensional space of
horizontal symmetric forms and a single positive ray. This fact is the
source of the uniqueness phenomenon in the next section.

\section{A positive Ermakov--Pinney dichotomy}\label{sec:pinney}

\subsection{Existence for every surreal potential}
For $q\in F$, consider the intrinsic scalar equations
\begin{equation}\label{eq:linearq}
 y''+qy=0
\end{equation}
and
\begin{equation}\label{eq:pinney}
 \rho''+q\rho=\rho^{-3},\qquad \rho>0.
\end{equation}
Primes always mean $\partial$. The sign on the right side is important:
we consider the positive Pinney constant $+1$.

\begin{theorem}[Positive Pinney existence and uniqueness]\label{thm:pinney}
For every $q\in\No$, equation \eqref{eq:pinney} has a positive surreal
solution. Exactly one of the following alternatives holds.
\begin{enumerate}
\item Equation \eqref{eq:linearq} has no nonzero solution in $F$. Then
\eqref{eq:pinney} has exactly one positive solution in $F$.
\item Equation \eqref{eq:linearq} has a real fundamental pair. Then the
positive solutions of \eqref{eq:pinney} form the classical positive
quadratic-form family, with two real parameters after determinant
normalization.
\end{enumerate}
In the first alternative, $y''+qy=f$ has a unique surreal solution for
every $f\in F$.
\end{theorem}
\begin{proof}
Let $A_q=\left(\begin{smallmatrix}0&1\\-q&0\end{smallmatrix}\right)$.
By \cref{thm:realmetrics}, it has a positive symmetric horizontal metric
$G$. The determinant identity for \eqref{eq:metric} gives
\[
 (\det G)'=-2\tr(A_q)\det G=0.
\]
Thus $\det G$ is a positive real constant. Scale $G$ by its inverse square
root to arrange $\det G=1$.

Write $G=\left(\begin{smallmatrix}a&b\\b&c\end{smallmatrix}\right)$.
The horizontal equation is exactly
\begin{equation}\label{eq:metric-components}
 a'=2qb,\qquad b'=qc-a,\qquad c'=-2b.
\end{equation}
Since $c>0$, put $\rho=\sqrt c$. The last equation gives $b=-\rho\rho'$;
the determinant condition gives $a=(\rho')^2+\rho^{-2}$. The middle
equation then becomes \eqref{eq:pinney}. Conversely, every positive
solution of \eqref{eq:pinney} produces the positive determinant-one metric
\begin{equation}\label{eq:PinneyMetric}
 G_\rho=
 \begin{pmatrix}
 (\rho')^2+\rho^{-2}&-\rho\rho'\\
 -\rho\rho'&\rho^2
 \end{pmatrix}.
\end{equation}
Positivity follows from its positive lower diagonal entry and determinant
one. This establishes a bijection between positive Pinney solutions and
positive determinant-one horizontal metrics.

In dimension two, the real normal form has either two zero phases or one
nonzero conjugate pair. If \eqref{eq:linearq} has no nonzero real solution,
it must be the second case. The positive metric cone is then a single
positive ray; determinant one selects one point, proving uniqueness.
In the zero-phase case, a real fundamental matrix exists and the metric
cone is the three-dimensional positive symmetric constant-matrix cone.
The determinant-one slice gives the asserted two-parameter family.
Finally, \cref{cor:inhomogeneous} gives a unique complex solution to every
forced equation in the nonzero-phase case; conjugating it shows that it is
real when $f$ is real.
\end{proof}

For completeness, the absence of an intermediate one-dimensional real
solution space can also be seen directly. If $u\in F^\times$ satisfies
\eqref{eq:linearq}, choose $B$ with $B'=u^{-2}$ and put $v=uB$. Then
$uv'-u'v=1$, so $u,v$ are an independent real fundamental pair.

\subsection{What is classical in the formula}
If $u,v\in F$ are a real fundamental pair with constant Wronskian
$W=uv'-u'v\ne0$, the second alternative is explicitly
\begin{equation}\label{eq:pinney-classical}
 \rho=\sqrt{a u^2+2buv+c v^2},
 \qquad a>0,\quad c>0,\quad ac-b^2=W^{-2},
 \quad a,b,c\in\R.
\end{equation}
This is the classical Pinney parametrization~\cite{Pinney}, now interpreted
in an ordered differential field. The new-field phenomenon is that the
first alternative is possible and collapses the positive family to one
point. The manuscript does not claim priority for the classical formula
or for the abstract Schwarz-closed-field existence principle underlying
it. Its metric proof exposes why positivity and the phase multiplicities
control uniqueness.
For recent existence and uniqueness results on the Hardy-field side,
see \cite[Theorems A and B]{ADHSecond}; no new solution of its named
Boshernitzan conjecture is being claimed here.

For example, $q=1$ lies in the first alternative, and its unique positive
Pinney solution is $\rho=1$. For $q=-1$, the linear equation has
$e^{\omega},e^{-\omega}$ as solutions, and the second alternative gives
many positive amplitudes. The existence of a unique positive solution of a
second-order equation is not a contradiction to ordinary initial-value
theory: intrinsic surreal differentiation has no arbitrary two-real-number
initial-data interpretation.

\subsection{A canonical symplectic phase for an oscillatory potential}
We call the first alternative \emph{intrinsically oscillatory}, meaning
only that the nonzero homogeneous oscillations are missing from the real
base field. Let $\rho$ be its unique positive Pinney solution, and put
$h=\rho^{-2}$. There is a unique primitive $B$ of $h$ whose coefficient of
the monomial $1$ is zero. Write
\[
 B=p+\eta,\qquad p\in\PP,\quad\eta\in\mm.
\]

\begin{theorem}[Canonical amplitude--phase reduction]\label{thm:canonicalpinney}
For an intrinsically oscillatory $q$, one has $p>0$. With
\[
 P_\rho=\begin{pmatrix}\rho&0\\\rho'&\rho^{-1}\end{pmatrix},
 \qquad R(\eta)=\exp(\eta J),
\]
the matrix $S_q=P_\rho R(\eta)\in\SL_2(F)$ satisfies
\begin{equation}\label{eq:canonical-gauge}
 S_q^{-1}A_qS_q-S_q^{-1}S_q'=p'J.
\end{equation}
The amplitude $\rho$, normalized primitive $B$, phase $p$, and this
specified amplitude--infinitesimal-rotation gauge are canonical. No
uniqueness among all possible symplectic gauges is asserted.
\end{theorem}
\begin{proof}
A direct multiplication using \eqref{eq:pinney} gives
\[
 P_\rho^{-1}A_qP_\rho-P_\rho^{-1}P_\rho'=\rho^{-2}J=hJ.
\]
The infinitesimal exponential $R(\eta)$ is a well-defined real rotation,
commutes with $J$, has determinant one, and satisfies
$R^{-1}R'=\eta'J$. This proves \eqref{eq:canonical-gauge}.
If $p=0$, then $B=\eta$ is finite and this same calculation reduces the
system to zero, supplying a real fundamental matrix, a contradiction.
Since $B'=h>0$, $B$ cannot be negative infinite by the $H$-field rule.
Therefore its nonzero purely infinite part is positive. The uniqueness
claims follow from \cref{thm:pinney} and the chosen constant-term
normalization of a primitive.
\end{proof}

The conserved quadratic expression is
\begin{equation}\label{eq:LewisInvariant}
 \mathcal E_\rho(y,y')=(\rho y'-\rho'y)^2+(y/\rho)^2.
\end{equation}
Its derivative vanishes formally under \eqref{eq:linearq}, by
\eqref{eq:PinneyMetric}. This expression is meaningful in an oscillatory
extension even when the only solution in $F$ is zero. It gives an exact
positive invariant, not an averaged energy.

\section{The oscillatory Airy equation: an exact Hahn example}\label{sec:airy}

\subsection{A formal oscillation with a genuine surreal amplitude}
Let $x=\omega$ and consider
\begin{equation}\label{eq:airy}
 y''+xy=0.
\end{equation}
Put $p=\tfrac23 x^{3/2}$, so $p'=x^{1/2}$. A formal oscillatory symbol
$X^p$ will mean a new invertible element with
$(X^p)'=ip'X^p$; it is \emph{not} a surcomplex number. Seek a solution
$aX^p$, where
\begin{equation}\label{eq:airyamp}
 a=x^{-1/4}\sum_{n\ge0}c_nx^{-3n/2},\qquad c_0=1.
\end{equation}
Substitution into \eqref{eq:airy} gives
\begin{equation}\label{eq:airyamp-eq}
 a''+2ix^{1/2}a'+\frac{i}{2x^{1/2}}a=0.
\end{equation}
Writing $\alpha_n=-\tfrac14-\tfrac32n$, coefficient comparison yields
\begin{equation}\label{eq:airyrec}
 c_{n+1}=\frac{\alpha_n(\alpha_n-1)}{3i(n+1)}c_n
       =-i\frac{(6n+1)(6n+5)}{48(n+1)}c_n.
\end{equation}
Equivalently,
\[
 c_n=(-i)^n\left(\frac34\right)^n
          \frac{(1/6)_n(5/6)_n}{n!}.
\]
The exponents in \eqref{eq:airyamp} decrease strictly to negative infinity,
so the displayed expression is a well-defined surcomplex Hahn sum,
regardless of factorial growth in the ordinary coefficients. Its termwise
derivative is justified by the strong property of the fixed derivation.
The recurrence proves \eqref{eq:airyamp-eq} exactly.

\subsection{A full phase gauge and its Wronskian}
Set
\[
 b=a'+ix^{1/2}a,\qquad
 P=\begin{pmatrix}a&\bar a\\b&\bar b\end{pmatrix}.
\]
The two columns satisfy $v'=A_xv-ip'v$ and its conjugate, where
$A_x=\left(\begin{smallmatrix}0&1\\-x&0\end{smallmatrix}\right)$. Thus
\[
 P'=A_xP-P\diag(ip',-ip').
\]
The determinant $W=a\bar b-\bar a b$ has derivative zero. Its leading term
is $-2i$, since $a\sim x^{-1/4}$ and $b\sim ix^{1/4}$. A differential
constant belongs to $\C$, so $W=-2i$ exactly. Hence $P$ is invertible and
\begin{equation}\label{eq:airyphases}
 \Specph(A_x)=\left\{\frac23\omega^{3/2},-\frac23\omega^{3/2}\right\}.
\end{equation}
It follows that \eqref{eq:airy} has no nonzero solution in $\K$, and that
$y''+\omega y=f$ has a unique surcomplex solution for every $f\in\K$.
This statement concerns intrinsic number solutions. It does not deny the
usual Airy functions of an ordinary real or complex variable.

\subsection{The unique positive nonlinear amplitude}
The unique positive solution of
\[
 \rho''+\omega\rho=\rho^{-3}
\]
is
\begin{equation}\label{eq:airyrho}
 \rho=\sqrt{a\bar a}.
\end{equation}
To verify the normalization without relying on a heuristic WKB argument,
form
\[
 C=\operatorname{Re}(vv^*)
 =\begin{pmatrix}|a|^2&\operatorname{Re}(a\bar b)\\
                  \operatorname{Re}(a\bar b)&|b|^2\end{pmatrix},
 \qquad v=\binom ab.
\]
The equation for $v$ implies $C'=A_xC+CA_x^{\mathsf T}$, because the scalar
imaginary phase cancels. Also
\[
 \det C=\bigl(\operatorname{Im}(a\bar b)\bigr)^2=1,
\]
as $W=-2i$. Therefore
\begin{equation}\label{eq:airymetric}
 G=C^{-1}=
 \begin{pmatrix}|b|^2&-\operatorname{Re}(a\bar b)\\
 -\operatorname{Re}(a\bar b)&|a|^2\end{pmatrix}
\end{equation}
is a positive determinant-one horizontal metric. The metric--Pinney
bijection proves \eqref{eq:airyrho}, and the nonzero phases in
\eqref{eq:airyphases} prove its uniqueness.

The expansion begins
\begin{equation}\label{eq:rhojet}
 \rho=\omega^{-1/4}\left(
 1-\frac5{64}\omega^{-3}
 +\frac{2285}{8192}\omega^{-6}
 -\frac{1690275}{524288}\omega^{-9}
 +\frac{10371973875}{134217728}\omega^{-12}
 +\cdots\right).
\end{equation}
These are coefficients of an exact Hahn series, not estimates for an
ordinary convergent power series. The normalized primitive of $\rho^{-2}$
has the form
\[
 B=\frac23\omega^{3/2}-\frac5{48}\omega^{-3/2}+\cdots,
\]
so its purely infinite part is exactly the positive phase in
\eqref{eq:airyphases}. The lower-order phase correction is finite and is
absorbed by the infinitesimal rotation in \cref{thm:canonicalpinney}.

\subsection{A residual certificate for truncation}
Let $a_N=\sum_{n=0}^N c_nx^{\alpha_n}$. The recurrence cancels every term
except the final second derivative, giving the exact identity
\begin{equation}\label{eq:airyres}
 a_N''+2ix^{1/2}a_N'+\frac{i}{2x^{1/2}}a_N
    =\alpha_N(\alpha_N-1)c_Nx^{\alpha_N-2}.
\end{equation}
The verification script checks this identity at $N=16$, as well as the
coefficient recurrence and the displayed amplitude coefficients. The
infinite-support proof is the recurrence plus Hahn summability; a finite
symbolic check by itself would not establish it.

\section{Polynomial and rational first integrals from phase resonance}\label{sec:invariants}

\subsection{Formal state variables and their total derivation}
Let $y_1,\ldots,y_n$ be algebraically independent formal variables and
extend the derivation to $\K[y_1,\ldots,y_n]$ by
\[
 \mathcal D_A f=\partial_{\mathrm{coeff}}f
   +\sum_j(Ay)_j\frac{\partial f}{\partial y_j}.
\]
A polynomial or rational \emph{first integral} means an element killed by
this derivation. This definition does not presuppose a nonzero solution in
$\K^n$. Under a phase gauge $y=Pz$, one has $\mathcal D_A z_j=ip_j'z_j$.

Let
\[
 \Lambda=\sum_{j=1}^n\Z p_j\subset\PP,\qquad
 L=\left\{m\in\Z^n:\sum_jm_jp_j=0\right\},\qquad
 S=L\cap\N^n.
\]
Write $r=\dim_{\Q}\Span_{\Q}\{p_1,\ldots,p_n\}$, so $\Lambda$ is a free
abelian group of rank $r$ and $L$ has rank $n-r$.

\begin{theorem}[Resonance algebra]\label{thm:invariants}
In phase coordinates, the polynomial first-integral algebra and rational
first-integral field are respectively
\begin{equation}\label{eq:invariantring}
 \ker\mathcal D_A\big|_{\K[z]}=\C[z^m:m\in S],
 \qquad
 \ker\mathcal D_A\big|_{\K(z)}=\C(z^m:m\in L).
\end{equation}
The polynomial algebra is finitely generated. The rational field is purely
transcendental over $\C$ of transcendence degree $n-r$.
\end{theorem}
\begin{proof}
For a polynomial $f=\sum_\alpha c_\alpha z^\alpha$, its derivative vanishes
if and only if
\[
 c_\alpha'+i\left(\sum_j\alpha_jp_j\right)'c_\alpha=0
\]
for each exponent. By \eqref{eq:phase-no-solution}, a nonzero coefficient
requires $\alpha\in S$ and $c_\alpha\in\C$. This proves the polynomial
formula.

For the rational statement, work in the Laurent polynomial ring and write
$f/g$ with coprime nonzero Laurent polynomials. If its derivative is zero,
then $f\mid\mathcal D_A f$ and $g\mid\mathcal D_A g$. The derivative does
not enlarge Laurent support. For each individual variable, the highest
and lowest Laurent degrees are additive under multiplication. Consequently
$\mathcal D_A f=hf$ and $\mathcal D_A g=hg$ for some $h\in\K$; the case
of a zero derivative is included by $h=0$.

Compare two coefficients of $f$. Their ratio has logarithmic derivative
$i(\beta\cdot p-\alpha\cdot p)'$, so their phases must agree and the
ratio is constant. Thus all monomials of $f$ have one phase, and after
factoring one coefficient and one monomial the remaining Laurent polynomial
has constant coefficients and exponents in $L$. The same holds for $g$.
The common value of $h$ forces their two factored phases to agree and their
factored coefficient ratio to be constant. Hence $f/g$ is a rational
function of the monomials with exponents in $L$, proving the second
formula. The reverse inclusion is immediate.

To prove finite generation of $S$, call a nonzero element indecomposable
when it is not a sum of two nonzero elements of $S$. Two distinct
indecomposables cannot be comparable in the coordinatewise order on
$\N^n$: if $u<v$, then $v-u\in S\setminus\{0\}$. Dickson's elementary
finite-antichain lemma for $\N^n$ therefore makes the indecomposables
finite. Induction on total degree writes every member of $S$ as a sum of
them. Finally, a $\Z$-basis of $L$ gives $n-r$ monomials which generate
the rational field and are algebraically independent: distinct integer
combinations of basis exponents give distinct Laurent monomials.
\end{proof}

One elementary proof of the finite-antichain lemma used here is by
induction on the number of coordinates: every infinite sequence of
nonnegative integers has an infinite nondecreasing subsequence, and
successively extracting such subsequences for finitely many coordinates
produces two coordinatewise comparable terms. No analytic finiteness
principle enters the argument.

\subsection{A coupled system with a singular invariant algebra}
Take $x=\omega$ and
\[
 P=\begin{pmatrix}1&x^{-1}&0\\0&1&x^{-2}\\0&0&1\end{pmatrix},
 \qquad D=\diag(i,i,-2i),\qquad A=P'P^{-1}+PDP^{-1}.
\]
Explicitly,
\begin{equation}\label{eq:3matrix}
 A=\begin{pmatrix}
 i&-x^{-2}&x^{-4}\\
 0&i&(-3ix-2)/x^3\\
 0&0&-2i
 \end{pmatrix}.
\end{equation}
Its phases are $\{\omega,\omega,-2\omega\}$, so it has no nonzero solution
in $\K^3$. Nevertheless, its formal polynomial first integrals are
nontrivial. Put $z=P^{-1}y$. The condition on an exponent is
$\alpha_1+\alpha_2=2\alpha_3$, and the invariant algebra is
\begin{equation}\label{eq:toricexample}
 \C[U,V,W]/(UW-V^2),\qquad
 U=z_1^2z_3,\quad V=z_1z_2z_3,\quad W=z_2^2z_3.
\end{equation}
Indeed these are the indecomposable nonnegative resonances; reduction by
the relation leaves distinct exponent monomials, so it generates the whole
kernel of the presentation map. The rational invariant field is
$\C(z_1/z_2,z_2^2z_3)$, of transcendence degree two. The cone in
\eqref{eq:toricexample} shows that complete reducibility of the differential
module does not force its polynomial invariant algebra to be a polynomial
ring. The accompanying script checks all three first-integral identities
in the original coupled coordinates.

\section{Finite oscillatory extensions and their Galois tori}\label{sec:galois}

\subsection{The finite phase lattice, not a class-sized choice of generators}
Let $A$ have phase multiset $\{p_1,\ldots,p_n\}$ and put
\[
 \Lambda=\Z p_1+\cdots+\Z p_n\subset\PP,
 \qquad r=\dim_{\Q}\Span_{\Q}\{p_1,\ldots,p_n\}.
\]
This is a finitely generated torsion-free abelian group, hence free of
rank $r$. Choose a $\Z$-basis $q_1,\ldots,q_r$ and form the Laurent
polynomial differential algebra
\begin{equation}\label{eq:PVring}
 R_\Lambda=\K[T_1^{\pm1},\ldots,T_r^{\pm1}],
 \qquad T_j'=i q_j'T_j.
\end{equation}
For $q=\sum_jm_jq_j$ write $X^q=\prod_jT_j^{m_j}$. Thus
$X^{q+s}=X^qX^s$ and $(X^q)'=iq'X^q$. Every expression in $R_\Lambda$
is a \emph{finite} sum. There is no Hahn summability question and no
attempt to define $e^{ip}$ in $\K$.

The general universal-exponential construction and its connection with
Picard--Vessiot theory are established theory~\cite[\S2.4]{ADHNorm}.
The point of the following proof is to identify exactly the finite
extension attached to the normal form, without invoking an existence
theorem for extensions of a proper-class field.

\begin{theorem}[The phase-lattice Picard--Vessiot algebra]\label{thm:PV}
The algebra $R_\Lambda$ is differentially simple, and the constants of
$E_\Lambda=\Frac(R_\Lambda)$ are exactly $\C$. If $P$ is a phase-normalizing
gauge for $A$, then
\[
 Y=P\diag(X^{p_1},\ldots,X^{p_n})
\]
is a fundamental matrix for $y'=Ay$. Its entries and inverse determinant
generate $R_\Lambda$ over $\K$. The differential Galois group is
\begin{equation}\label{eq:galoistorus}
 \operatorname{Gal}_\partial(E_\Lambda/\K)
   =\Hom(\Lambda,\C^\times)\simeq(\C^\times)^r,
\end{equation}
with its algebraic-group structure the split torus $\mathbb G_m^r$.
In particular, $\operatorname{trdeg}_{\K}E_\Lambda=r$.
\end{theorem}
\begin{proof}
Suppose $J\ne0$ is a differential ideal of $R_\Lambda$. Choose an element
of $J$ with a minimum number of nonzero monomials. Multiply it by a
Laurent monomial and a nonzero coefficient to arrange that its constant
term is $1$. Write the result as $f=1+\sum_{q\ne0}c_qX^q$.
Differentiation kills its constant term and cannot introduce new
monomials. If $f'\ne0$, it contradicts minimality. If $f'=0$, each
remaining coefficient satisfies
\[
 c_q'/c_q=-i q',
\]
which is impossible for nonzero $q\in\PP$ by the rank-one phase
criterion. Therefore $f=1$ and $J=R_\Lambda$.

To check the constant field directly, let $(f/g)'=0$ with coprime
nonzero Laurent polynomials $f,g$. Then $f'g=fg'$, so $f\mid f'$ and
$g\mid g'$. Differentiation preserves or reduces the range of powers of
each $T_j$. Since the largest and smallest powers in a product add,
the quotients must lie in $\K$: $f'=hf$ and $g'=hg$ for a common
$h\in\K$. Two different monomials $c_qX^q,c_sX^s$ in $f$ would give
\[
 (c_q/c_s)'/(c_q/c_s)=-i(q-s)',
\]
contradicting $q-s\ne0$. Thus $f$ and $g$ are each monomials.
Their quotient can have derivative zero only if their exponents coincide
and their coefficient ratio lies in $\C$. This proves
$C_{E_\Lambda}=\C$.

The normal-form equation $P'=AP-PD_p$ immediately gives $Y'=AY$.
Conversely, the entries of $P^{-1}Y$ supply every $X^{p_j}$, and the
inverse determinant supplies their product inverse. Multiplying this by
all the factors except one supplies each $(X^{p_j})^{-1}$. These
characters generate exactly the group algebra of $\Lambda$, namely
$R_\Lambda$.

A differential $\K$-automorphism must take $X^{q_j}$ to
$c_jX^{q_j}$ with $c_j\in\C^\times$, because the ratio of two solutions
of the same rank-one equation is constant. Every choice of these $r$
constants defines an automorphism of the Laurent algebra and its fraction
field. This proves \eqref{eq:galoistorus}; the description as a Laurent
rational field also proves the transcendence-degree statement.
\end{proof}

Equivalently, in column coordinates the Galois group is the subgroup of
$(\C^\times)^n$ cut out by
\[
 \prod_{j=1}^n z_j^{m_j}=1
 \qquad\text{for every }m\in\Z^n\text{ with }\sum_jm_jp_j=0.
\]
The integer relation lattice is saturated: if a nonzero integer multiple
of $\sum_jm_jp_j$ vanishes, that sum already vanishes. This explains why
the resulting torus is connected and has no extra finite component group.
The rank is a \emph{rational} rank. The phases $\omega$ and
$\sqrt2\,\omega$ give a two-dimensional torus, although their real span
has dimension one. The phases $\omega$ and $2\omega$ give only one
dimension. The three-phase example \eqref{eq:3matrix} also gives rank one.

\subsection{Real structure and the scope of the Galois conclusion}
For a real system, its phase lattice is stable under negation.
Complex conjugation extends to an involution of $R_\Lambda$ by
\[
 \overline{cX^q}=\bar cX^{-q}.
\]
It commutes with differentiation. The combinations
$(X^q+X^{-q})/2$ and $(X^q-X^{-q})/(2i)$ are fixed and supply the
formal cosine and sine entries of a real fundamental matrix in the
real block normal form. An automorphism commuting with this involution
has $\chi(q)\overline{\chi(q)}=1$. Thus the real form of the torus is
\[
 \bigl(\operatorname{Res}^{1}_{\C/\R}\mathbb G_m\bigr)^r,
\]
whose real points are $(S^1)^r$. This is an algebraic-group assertion,
not a construction of trigonometric functions on all surreals.

The statement that all these finite Galois groups are tori concerns the
very large base field $\K$, with its factorization and inhomogeneous
surjectivity properties. It says nothing comparable for an equation over
its original small rational-function field. In particular, the Airy
calculation is not a claim that the usual Airy differential Galois group
over $\C(x)$ is a torus. Enlarging the base to $\K$ has already supplied
the nonoscillatory gauge entries.

When $\Lambda\ne0$, no differential embedding $E_\Lambda\to\K$ fixing
$\K$ can exist: the image of $X^q$ for nonzero $q$ would solve an
obstructed rank-one equation in $\K$. The extension does not erase the
original obstruction; it represents it explicitly.

\section{Foundations, proof audit, and reproducibility}\label{sec:audit}

\subsection{Handling the proper class without a class-sized basis theorem}
The full surreal field is a proper class. One possible formal setting
for the present exposition is class theory with global choice, reading
finite-dimensional algebra as a scheme about class operations on
set-coded finite tuples. No proper class is used as an element of a set,
and no proof needs a basis of $\PP$ as a class-sized vector space.
Every phase lattice above is finitely generated, and every Laurent
polynomial has finitely many terms.

There is also a set-sized interpretation of each finite-data argument.
Starting from a set of given surreal coefficients and $\R\cup\{\omega\}$,
form a countable chain of set-sized subfields inside $F$. At each stage
adjoin the real algebraic roots, derivatives, real exponentials, purely
infinite parts, primitives, and finite-angle trigonometric values needed
for the previous stage. Also adjoin witnesses for all finite-order
factorizations and inhomogeneous first-order equations whose coefficients
belong to that stage. There are only set many such requests at each
stage, and every request has a witness in $F$ or $F[i]$ by the established
inputs. Global choice permits the simultaneous selections. Taking field
closures and a countable union still yields a set.

The resulting real differential subfield $F_0\subset F$ is real closed,
contains $\R$, inherits the $H$-field rule, and has the designated
primitive and linear-equation witnesses. Its complexification has
constant field exactly $\C$. All finite tuples of its elements occur
at some stage, so their required witnesses occur later. The elementary
proofs of the finite-primitive and rank-one lemmas then hold internally,
as do the finite matrix arguments. Explicit Hahn series that are to be
used can be included in the initial set. One may therefore interpret
\cref{thm:PV} as an ordinary set-sized Picard--Vessiot statement over a
suitable witness-closed $F_0[i]$, as well as through its explicit finite
class-algebra formulas over $\K$.

This construction does \emph{not} say that an arbitrary fixed Hahn
workspace already has every needed differential primitive or solution.
Nor does it give effective presentations, birthday bounds, or a minimal
such subfield. It is a size justification for the proof, not an algorithm.

\subsection{The dependency chain and its load-bearing steps}
The proof of the principal theorem has a short logical structure even
though establishing its hypotheses is deep:
\begin{enumerate}
\item The Berarducci--Mantova and Aschenbrenner--van den Dries--van der
Hoeven results give the $H$-field rule, primitives, scalar exponential
compatibility, linear splitting, and inhomogeneous first-order
surjectivity.
\item The $H$-field rule makes $\II$ an order-convex $\OO$-module.
The polar logarithm gives the exact scalar phase obstruction.
\item First-order factorization and surjectivity give the diagonal
phase normal form. Horizontal Hermitian forms make that normal form
unitary when the original coefficient is skew-Hermitian.
\item A unitary matrix is bounded, so every entry of its connection
term belongs to $\II_{\C}$. The algebraic Weyl estimate puts the changes
of the ordered eigenvalues in $\II$ too. Applying $\Psi$ removes exactly
these changes and proves \cref{thm:main}.
\end{enumerate}
Only the first step is imported at this level of generality. The
intermediate finite-dimensional, order, and gauge arguments are proved
in the manuscript. In particular, neither an asymptotic adiabatic theorem
nor an unproved claim that diagonalization commutes with differentiation
is hidden in the fourth step.

The main calculation also exposes two essential boundaries. Removing
Hermitian symmetry permits the counterexample in \cref{ex:nonnormal}.
Replacing the finite-primitive ideal by all infinitesimals permits the
scalar counterexample $\omega^{-1}$ in \cref{prop:sharp}. Neither
hypothesis is decorative.

\subsection{What was checked computationally}
The archive contains \texttt{code/verify\_examples.py}, requiring Python
and SymPy. Its recorded execution with SymPy~1.14.0 passed \textbf{62}
exact checks. These comprise the nonnormal isospectral counterexample;
the rational unitary coupled example; a positive horizontal metric with
an unbounded gauge; the three-phase normal form and its polynomial
invariants; the noncommuting matrix recurrence and its unitary jet; the
Airy coefficient recurrence, amplitude jet, and exact truncation
residual; and the symbolic Pinney metric and symplectic gauge identities.

The deliberate counterexamples are checked as counterexamples: the
nonnormal coefficient has eigenvalues $\pm i$ but a fundamental matrix
in the base field, and the noncommuting exponential defect is required
to be \emph{nonzero}. The report is in
\texttt{data/verification.json}. Re-running the script regenerates it.
The checks use exact rational and symbolic arithmetic, not numerical
agreement at selected real sample points.

These tests do not verify the arbitrary-support Hahn arguments, the
published model-theoretic inputs, the existence of a gauge for an
arbitrary surreal matrix, or the general theorems. The article is not a
Lean formalization and no proof-assistant verification is claimed.
Its strongest claims require independent mathematical review.

\subsection{Limitations and questions left genuinely open}
The main novelty assessment is provisional. The targeted search located
the spectral and differential-algebra ingredients but not the exact
Hermitian spectral--integral formula, its sharp entrywise perturbation
criterion, or the normalized non-Abelian integration statement in this
form. This is evidence for treating them as candidate new results, not
proof of historical priority. In contrast, the rank-one obstruction,
universal exponential algebra, classical Pinney formula, and general
amplitude--phase methods have explicit antecedents.

All results depend on the specified intrinsic derivation. They are not
claims about every surreal derivation, arbitrary class-function
derivatives, or the independent complex-coordinate calculus used in other
repository reports. Infinite-dimensional matrices, topological initial
value problems on $\No$, Stokes data over smaller analytic fields, and a
global surcomplex exponential are outside the assertions.

The phase map and the normalizing gauge are semantic operations. Even
for finite matrices, arbitrary surreal input has no assumed finite
computable encoding. A useful further problem is to identify restricted
support-certified input languages for which the eigenvalue primitives,
their purely infinite parts, and a normalizing unitary gauge admit
certified effective computation. The exact recurrence examples here are
instances, not a solution of that general effectivity problem.

\section{Conclusion}
For intrinsic surcomplex differential systems, the obstruction to
integration is not exhausted by scalar examples, but it is still exactly
controlled by scalar phases after the necessary finite-dimensional
reduction. Hermitian symmetry makes this classification unusually
explicit: integrate the algebraic eigenvalues, retain their purely
infinite parts, and obtain the complete differential unitary class.
The connection term is neither ignored nor approximated; it belongs to
the precise finite-primitive ideal that the invariant annihilates.

This yields exact, gap-free perturbation invariance and a non-Abelian
integration theorem. The same phase decomposition classifies all positive
horizontal metrics, resolves the real two-dimensional amplitude problem,
and identifies the finite resonances and oscillatory extension needed
for any coupled system. The results therefore connect surreal order,
normal-form truncation, finite spectral algebra, and differential Galois
theory in a way that is not merely a replacement of ordinary scalar
symbols by surreal ones.

\begingroup\small\raggedright
\begin{thebibliography}{99}
\setlength{\itemsep}{.25em}

\bibitem{BM}
A.~Berarducci and V.~Mantova,
\emph{Surreal numbers, derivations and transseries},
Journal of the European Mathematical Society \textbf{20} (2018),
no.~2, 339--390.
\href{https://doi.org/10.4171/JEMS/769}{doi:10.4171/JEMS/769}.
Preprint: \href{https://arxiv.org/abs/1503.00315}{arXiv:1503.00315}.

\bibitem{ADHSur}
M.~Aschenbrenner, L.~van den Dries, and J.~van der Hoeven,
\emph{The surreal numbers as a universal $H$-field},
Journal of the European Mathematical Society \textbf{21} (2019),
no.~4, 1179--1199.
\href{https://doi.org/10.4171/JEMS/858}{doi:10.4171/JEMS/858}.
Preprint: \href{https://arxiv.org/abs/1512.02267}{arXiv:1512.02267}.

\bibitem{ADHNorm}
M.~Aschenbrenner, L.~van den Dries, and J.~van der Hoeven,
\emph{Normalizing Asymptotic Differential Equations},
\href{https://arxiv.org/abs/2403.19732v5}{arXiv:2403.19732v5},
April~24, 2026.
Especially \S2.4 and Lemma~2.4.19; the earlier versions had the title
\emph{A Normalization Theorem in Asymptotic Differential Algebra}.

\bibitem{ADHSecond}
M.~Aschenbrenner, L.~van den Dries, and J.~van der Hoeven,
\emph{Revisiting Second-Order Linear Differential Equations over Hardy Fields},
\href{https://arxiv.org/abs/2603.02013v1}{arXiv:2603.02013v1},
March~2, 2026. Especially Theorems~A and~B and the amplitude--phase
formulations in the introduction.

\bibitem{Pinney}
E.~Pinney,
\emph{The nonlinear differential equation $y''+p(x)y+cy^{-3}=0$},
Proceedings of the American Mathematical Society \textbf{1} (1950),
no.~5, 681.
\href{https://doi.org/10.1090/S0002-9939-1950-0037979-4}
{doi:10.1090/S0002-9939-1950-0037979-4}.
The sign of the parameter is convention dependent; the equation used
in this article is $\rho''+q\rho=\rho^{-3}$.

\bibitem{RepoMap}
\texttt{VladimirReshetnikov/Surreal},
\emph{The surreal and surcomplex reports}, documentation map,
\texttt{docs/README.md}, repository revision \texttt{aa846271b4dc},
accessed September~21, 2026.
\href{https://github.com/VladimirReshetnikov/Surreal/blob/aa846271b4dcae2c055b216126a87210292ec19b/docs/README.md}
{Pinned repository source}.

\bibitem{RepoDE}
\texttt{VladimirReshetnikov/Surreal},
\emph{Differential Algebra and Differential Equations over the Surreal
and Surcomplex Numbers}, report inventory,
\texttt{docs/surcomplex/differential-equations/README.md},
repository revision \texttt{aa846271b4dc}, accessed September~21, 2026.
\href{https://github.com/VladimirReshetnikov/Surreal/blob/aa846271b4dcae2c055b216126a87210292ec19b/docs/surcomplex/differential-equations/README.md}
{Pinned repository inventory}.

\bibitem{RepoSpectral}
\texttt{VladimirReshetnikov/Surreal},
\emph{Surcomplex Spectral Theory: Singular Values, Infinitesimal Rank,
and Multiscale Stability}, opening portions of the source manuscript,
\texttt{docs/surcomplex/spectral-theory/article.tex},
repository revision \texttt{aa846271b4dc}, accessed September~21, 2026.
\href{https://github.com/VladimirReshetnikov/Surreal/blob/aa846271b4dcae2c055b216126a87210292ec19b/docs/surcomplex/spectral-theory/article.tex}
{Pinned repository source}.

\end{thebibliography}
\endgroup
\end{document}
